% Flavor catalog per-process draft.
% process_id: L003
% family: charged_lepton
% owner_agent_id: PKA-L003
% writer_agent_id: null
% checker_agent_id: null
% opus_signoff_id: null
% Canonical metadata sidecar: L003.yaml
% Status is controlled by the YAML sidecar status_history, not this comment.

\section{L003: \texorpdfstring{\(\mu-e\)}{mu-e} Conversion in Aluminum}

\subsection*{Process}
\textbf{Standard notation:} \(R_{\mu e}^{\rm Al} =
\Gamma(\mu^-{\rm Al}\to e^-{\rm Al})/\Gamma_{\rm capture}(\mu^-{\rm Al})\).
This entry covers coherent neutrinoless muon-to-electron conversion on an
aluminum target, the design target for Mu2e and COMET.  It is a
charged-lepton-flavor-violating nuclear process, not part of the quark
\(\Delta F=2\) scan lane.

\subsection*{PDG / equivalent value(s)}
There is no published direct aluminum-target limit in the sources snapshotted
for this PKA pass.  The current published benchmark for coherent conversion is
the SINDRUM II gold result,
\[
  R_{\mu e}^{\rm Au} < 7\times10^{-13}\quad (90\%\ {\rm C.L.}),
\]
recorded in \texttt{L003.yaml:current\_world\_limit} and snapshotted at
\texttt{flavor\_catalog/references/L003/sindrumii\_2006\_gold\_juser.txt}.
For aluminum, the catalog value is therefore prospective.  The Mu2e
full-exposure projection gives single-event sensitivity
\(2.8\times10^{-17}\) and an expected \(90\%\) C.L. upper limit
\(6.7\times10^{-17}\).  Mu2e Run I separately quotes an expected no-signal
limit \(R_{\mu e}<6.2\times10^{-16}\) and discovery sensitivity
\(1.2\times10^{-15}\).  COMET Phase-I quotes \(3.1\times10^{-15}\)
single-event sensitivity and a \(7\times10^{-15}\) expected upper limit, while
the COMET Phase-II planning snapshot quotes \(2.6\times10^{-17}\).

\subsection*{Relevance to the RS / anarchic-flavor pipeline}
In a warped lepton-flavor extension, coherent conversion probes dipole,
scalar, and vector lepton-quark operators.  KK gauge exchange, \(Z\)-like
flavor violation, Higgs-mediated scalar currents, and lepton-sector dipoles can
all feed the same nuclear final state.  The observable is especially useful
because it is often more sensitive to contact interactions than
\(\mu\to e\gamma\), so it tests parts of the lepton extension that the current
repo dipole checker cannot see.

\subsection*{Post-2008 developments}
Relative to the arXiv:0804.1954 RS-flavor baseline, the main development is
experimental: aluminum conversion became a flagship next-generation CLFV
target.  Mu2e published staged Run I sensitivity projections, and Mu2e-II
records a proposed PIP-II upgrade aiming for at least another order of
magnitude beyond Mu2e.  COMET published a Phase-I technical design and a
staged Phase-II strategy.  These developments move the process from a
SINDRUM-II-era contextual bound to a near-term discovery channel with projected
sensitivity roughly four orders below the gold benchmark.

\subsection*{Constraint validity and model dependence}
The experimental signature is robust: a monoenergetic conversion electron from
a stopped negative muon with no neutrinos.  The mapping to RS parameters is
model-dependent.  It requires a lepton-quark operator basis, target-dependent
overlap integrals and capture rates, and a convention for combining dipole,
scalar, and vector contributions.  A gold limit cannot be treated as an
aluminum limit without a nuclear/operator translation.

\subsection*{Code coverage in this repo}
\textbf{NO.}  The required greps over \texttt{quarkConstraints/}, \texttt{qcd/},
\texttt{flavorConstraints/}, \texttt{neutrinos/}, \texttt{yukawa/},
\texttt{warpConfig/}, \texttt{solvers/}, \texttt{scanParams/}, and
\texttt{tests/} found no Mu2e, COMET, SINDRUM, \(R_{\mu e}\), or
\(\mu-e\)-conversion implementation.  The only adjacent LFV code is the
\(\mu\to e\gamma\) dipole checker at
\texttt{flavorConstraints/muToEGamma.py}:75 and its scan call at
\texttt{scanParams/scan.py}:524.

\subsection*{Implementation difficulty}
\textbf{HIGH.}  A production implementation needs new lepton-quark Wilson
coefficients, target-dependent conversion formulae for aluminum, nuclear
overlap/capture inputs, and likely EFT running or matching shared with
\(\mu\to e\gamma\) and \(\mu\to3e\).  This is beyond adding a new numerical
limit to the existing dipole-only LFV path.

\subsection*{Key references}
Process-local source keys before bibliography consolidation:
\texttt{SINDRUMII2006\_GoldMuE}, \texttt{Mu2e2016\_FullProjection},
\texttt{Mu2eRunI2023}, \texttt{COMETPhaseITDR2020},
\texttt{COMETStrategy2018}, \texttt{Mu2eII2022}, and \texttt{CFW2008}.
